% A3 — draft v0.1 (18/07/2026). Target: BRACIS (LNCS) → versão estendida em
% Applied Soft Computing / Neurocomputing. Class: article (converter depois).
\documentclass[11pt]{article}
\usepackage[T1]{fontenc}
\usepackage[utf8]{inputenc}
\usepackage[english]{babel}
\usepackage{booktabs}
\usepackage{amsmath}
\usepackage[round]{natbib}
\usepackage[margin=2.8cm]{geometry}
\usepackage{microtype}
\usepackage{xcolor}
\newcommand{\todo}[1]{\textcolor{red}{[TODO: #1]}}

\title{Label-free cold start for active learning beats a supervised
evolutionary envelope\\ --- if the envelope is honest: a circularity audit}

\author{Gilsiley Henrique Darú$^{1}$ \and Gustavo Valentim Loch$^{1}$\\[2pt]
\normalsize $^{1}$Graduate Program in Numerical Methods in Engineering (PPGMNE),\\
\normalsize Federal University of Paraná (UFPR), Curitiba, Brazil}
\date{Draft v0.1 --- \today}

\begin{document}
\maketitle

\begin{abstract}
The first labeled set $L_0$ decides how an active learning loop starts ---
and in short-text classification with hundreds of imbalanced classes, it
decides a lot: we measure a 6.4-point accuracy spread at $|L_0|=100$
caused by the random draw alone (47 sizes $\times$ 30 repetitions,
independently re-executed with $\le 0.7$-point divergence). To estimate
how much deliberate composition can yield, a genetic algorithm optimizes
$L_0$ composition under four objectives, revealing a wide envelope
(up to 25.9 points between best and worst individuals at $|L_0|=100$)
that narrows with size. We then show that DRI-SL --- a deterministic,
label-free, linear-cost heuristic combining semantic density
(embedding-cluster representativeness) with lexical variety ---
\emph{outperforms the best individual found by the supervised optimizer}
at every size from 100 to 5{,}000 (e.g., 67.4\% vs.\ 62.2\% accuracy at
$|L_0|=1{,}000$). Finally, a methodological audit shows the standard
protocol of reporting evolutionary envelopes on the same partition used as
fitness inflates them --- by 6.3 points of macro-F1 at $|L_0|=500$ in our
re-execution with an anti-circularity protocol --- which means published
evolutionary bounds for subset selection are optimistic and our
comparison is conservative. Composition, not just quantity, governs the
cold-start regime; it can be captured without any labels; and envelopes
must be reported on untouched partitions.
\end{abstract}

\noindent\textbf{Keywords:} active learning; cold start; instance
selection; genetic algorithms; evaluation methodology; short text.

\section{Introduction}
\label{sec:intro}

Pool-based active learning assumes a model good enough to guide selection
--- which presupposes an initial labeled set $L_0$. In the cold-start
regime the loop's later efficiency is hostage to this first draw
\citep{Bayer2024,Attenberg2010}. Three questions structure this paper:
(RQ1) how much does the \emph{composition} of $L_0$ matter, at fixed size?
(RQ2) how much headroom does deliberate composition have over random
draws --- an envelope we estimate with a supervised evolutionary
optimizer? (RQ3) can a label-free constructive heuristic capture that
headroom? A fourth, methodological question emerged from answering the
second: (RQ4) are evolutionary envelopes, as usually reported, even
honest?

We answer on a hard instance of the problem: 250k short product
descriptions (4--50 characters, abbreviations) in Portuguese, 621
imbalanced classes, where class coverage --- not instance count ---
dominates early performance.

\section{Related work}
\label{sec:related}

Cold-start-aware AL uses pretrained representations to pick informative
seeds: ALPS exploits language-model surprisal \citep{Yuan2020}; BERT-era
studies adapt AL to fine-tuning \citep{EinDor2020,Griesshaber2020}; DEUCE
combines textual and class-prediction diversity \citep{Guo2025Deuce};
LLM-based selectors push the idea further \citep{Bayer2024}.
Density/cluster-based selection is classical
\citep{Settles2008,Nguyen2004,Dasgupta2008}, as are batch-diversity
concerns \citep{Hoi2006}. Evolutionary algorithms
\citep{Holland1975,Goldberg1989} are widely used for instance and subset
selection; we are not aware of prior work auditing the \emph{evaluation
circularity} of their reported envelopes in this setting --- fitness
measured and reported on the same partition --- although the underlying
overfitting-to-selection phenomenon is well known in model selection.

\section{Setting}
\label{sec:setting}

Classifier: PVBin, a per-class prototype model over binary term vectors,
cheap enough for tens of thousands of retrainings
\citep{Daru2024Dissertacao}. Data: deduplicated short product
descriptions, 621 classes, heavy long tail; evaluation on a fixed
untouched test set (20{,}000 instances in the re-execution protocol).
Metrics: global accuracy and macro-F1, the latter primary under
imbalance \citep{Sokolova2009}.

\section{RQ1: sensitivity to the draw}
\label{sec:sens}

Across 47 sizes (10 to 200{,}000) with 30 random repetitions each:
mean accuracy grows from 6.7\% ($|L_0|=10$) through 24.7\% (100), 55.9\%
(1{,}000) and 76.9\% (10{,}000), saturating near 89.1\% (200{,}000);
macro-F1 lags (0.4\% $\to$ 3.4\% $\to$ 17.8\% $\to$ 43.3\% $\to$ 70.4\%),
reflecting the cost of covering minority classes. At $|L_0|=100$,
accuracy across repetitions spans \textbf{21.7--28.1\%} --- a
\textbf{6.4-point} spread caused exclusively by the draw. The mechanism
is coverage: at $|L_0|=10$ a typical draw covers 9 of 621 classes; at
1{,}000, about 255; macro-F1 tracks the coverage curve closely. Good and
bad $L_0$ of the same size \emph{exist}; the question becomes whether
they can be found without labels.

\paragraph{Independent re-execution.} The full study was re-executed in a
clean-room reimplementation with a stricter protocol (deduplication
\emph{before} partitioning; fixed test set), on a reduced grid (15
log-spaced sizes $\times$ 10 repetitions). Agreement is within 0.7 points
at every size, always in the expected direction (the deduplicated
protocol is slightly harder), with implementation fidelity additionally
verified by exact equality of score matrices.

\section{RQ2: the evolutionary envelope}
\label{sec:ga}

A genetic algorithm over subset membership (100 generations; the full
configuration ships in the released artifacts) optimizes $L_0$
composition for four objectives (max/min $\times$ accuracy/macro-F1) at ten sizes.
The envelope is wide where it matters: at $|L_0|=100$, best-minus-worst
final individuals span \textbf{25.9 accuracy points}; evolutionary gain
over the first generation (e.g., $+10.1$ points at 100, $+11.7$ at 50)
confirms the optimizer exploits real structure, not noise. The envelope
narrows monotonically with size (at 30{,}000: $+0.8$), consistent with
RQ1: composition is a small-regime lever.

\section{RQ3: a label-free heuristic vs.\ the supervised optimizer}
\label{sec:drisl}

DRI-SL (Density-Representative Instance Selection with Lexical variety)
is deterministic given the embedding model and uses no labels: embed
the pool with dense text embeddings, cluster ($k$-means), allocate the budget proportionally to
cluster mass, and fill each cluster's quota greedily preferring
instances that add unseen lexical tokens --- linear cost after
clustering. Table~\ref{tab:drisl} compares it with the GA envelope.

\begin{table}[htb]
\centering
\caption{Cold start: DRI-SL vs.\ GA envelope (final generation), PVBin.
GA values inherit the original protocol (fitness partition = report
partition) and are therefore optimistic --- see
Section~\ref{sec:audit}.}
\label{tab:drisl}
\small
\begin{tabular}{r rr rr rr}
\toprule
 & \multicolumn{2}{c}{\textbf{DRI-SL}} & \multicolumn{2}{c}{\textbf{GA (best)}}
 & \multicolumn{2}{c}{\textbf{GA (worst)}} \\
$|L_0|$ & Acc & M-F1 & Acc & M-F1 & Acc & M-F1 \\
\midrule
100    & \textbf{41.2} & \textbf{6.9}  & 38.8 & 6.5  & 5.8  & 1.8 \\
500    & \textbf{59.2} & \textbf{18.5} & 56.0 & 18.0 & 36.2 & 6.5 \\
1{,}000 & \textbf{67.4} & \textbf{25.8} & 62.2 & 24.7 & 49.3 & 12.6 \\
2{,}500 & \textbf{73.4} & \textbf{36.5} & 69.3 & 32.6 & 63.9 & 24.1 \\
5{,}000 & \textbf{76.9} & \textbf{44.1} & 74.1 & 40.1 & 70.8 & 31.9 \\
\bottomrule
\end{tabular}
\end{table}

A constructive heuristic with zero label access beats an optimizer that
consumed thousands of supervised evaluations, at every size --- strong
evidence that semantic density plus lexical variety capture precisely the
structure that makes an $L_0$ informative (class coverage through cluster
coverage; feature coverage through lexical novelty). At 100 the margin
over the GA's best is 2.5 points; over the GA's \emph{worst}, 35 points
--- the cost of a bad draw that DRI-SL eliminates deterministically.

\section{RQ4: the circularity audit}
\label{sec:audit}

The GA numbers above inherit the field's standard protocol: fitness
evaluated on a partition, and the \emph{same} partition reported as the
result. Our re-execution applies an anti-circularity protocol --- fitness
on a selection partition, final individuals \emph{re-evaluated once} on
the untouched test set. The mechanism reproduces (best individual gains
$+5.2$ accuracy points at $|L_0|=50$), but the honest re-evaluation
reveals the inflation: for macro-F1 maximization at $|L_0|=500$, final
fitness reads 19.4\% while the same individual measures \textbf{13.1\%}
on untouched data --- \textbf{6.3 points of overfitting to the fitness
partition} (for accuracy, $\approx 1$ point). Two consequences: (i)
published evolutionary envelopes for subset selection should be read as
optimistic bounds unless the protocol states otherwise; (ii) our central
comparison is \emph{conservative} --- DRI-SL beats an inflated opponent.

\section{Conclusion}
\label{sec:conclusion}

Composition of the initial labeled set is a first-order factor in
cold-start active learning for long-tailed short text (6.4-point spread
from the draw alone at $|L_0|=100$); a supervised evolutionary optimizer
confirms wide headroom; a deterministic label-free heuristic captures
that headroom, beating the optimizer's best individual up to
$|L_0|=5{,}000$; and evolutionary envelopes must be reported on untouched
partitions --- ours was inflated by up to 6.3 macro-F1 points before
correction. In the continuous (post-cold-start) regime, companion work
shows the lexical-variety component should be dropped in favor of plain
prediction-stratified sampling; the two results together delimit where
each ingredient earns its keep. All artifacts are released.

\paragraph{Artifacts.} \texttt{github.com/GHDaru/activelearning}
(\texttt{experiments/p1}, \texttt{p2}).

\bibliographystyle{plainnat}
\bibliography{refs}

\end{document}
